% ##############################################################################
\section{Introduction}
\label{sec:introduction}

A multi-tenant product makes promises. Every route, button, module, and sidebar entry says something to the user. It says: this belongs to you; this will work; this is part of your world.

Those promises are easy to break. The backend may enforce authorization correctly, while the frontend carries a separate permission model made from route guards, role checks, feature flags, and tenant-specific conditionals. The two models rarely fail in one dramatic incident. They drift one small decision at a time. A restricted tenant sees a billing module. A support engineer clicks into a page that was never enabled for that role. A module fetches configuration, discovers the user cannot use it, and renders a blank shell.

That is not only ugly product behavior. It is authorization drift. The UI is now describing a world the backend will not honor.

We adopt a blunt rule.

\begin{quote}
\textbf{If the UI can lie about access, it is part of the security boundary.}
\end{quote}

\uiPolicy{} is the discipline we use to make that boundary honest. The backend remains the enforcement point. The UI does not become trusted security code. Instead, the UI is rendered from a server-materialized contract that comes from the same policy inputs used for authorization. The frontend stops guessing.

The core primitive is \uiData{}: a tenant- and role-scoped UI contract containing enabled modules, layout fragments, and the effective support context. The web tier consumes that contract to compose navigation and module layouts. It also builds the client-side ability model from a server-side profile, then exports packed rules for user-experience decisions. Backend APIs continue to authorize every request.

\paragraph{Concrete failure.}
During onboarding for a restricted tenant, a frontend conditional for a billing module was not updated with the backend entitlement record. The sidebar showed the billing link. The backend returned HTTP 403. Three support escalations later, the discrepancy traced back to a stale if-block. Under \uiPolicy{}, the sidebar is rendered from the server-materialized module list; there is no standalone billing conditional to forget.

\paragraph{Contributions.}
This paper makes four contributions.

\begin{itemize}
  \item It names the representation gap between backend enforcement and frontend rendering, then states a Truthful UX invariant for policy-gated surfaces.
  \item It describes \uiData{}, a production pattern where tenant modules and layout fragments are materialized as a runtime UI contract instead of being scattered through frontend code.
  \item It shows how server-built CASL abilities, module layout checks, and backend enforcement keep the client useful without making it trusted.
  \item It presents a view-as support workflow with immutable actor attribution, server-enforced allow-lists, signed effective-context headers, and searchable audit logs.
\end{itemize}

The implementation uses Next.js, Cognito, AWS Lambda, DynamoDB, iron-session, and CASL \cite{aws_cognito_user_pools,aws_cognito_groups,aws_lambda,aws_dynamodb,iron_session,casl_docs}. The pattern itself is not tied to those choices.
